%====================================================
% PART B
% NIST Cybersecurity Framework 2.0
%====================================================

\chapter{NIST Cybersecurity Framework 2.0}

This chapter examines the NIST Cybersecurity Framework (CSF) Version 2.0, published by the National Institute of Standards and Technology in February 2024. It analyses the six core functions introduced in the updated framework, the four implementation tiers, the use of organizational profiles for gap analysis, and applies the framework to a hospital ransomware incident response scenario.

%====================================================
\section{Core Functions}
%====================================================

\subsection{Govern}

\subsubsection{Primary Purpose}

\subsubsection{Key Categories}

% Category 1:

% Category 2:

\subsubsection{Implementation Example}

\subsubsection{Interaction with Other Functions}

%---------------------------------------

\subsection{Identify}

\subsubsection{Primary Purpose}

\subsubsection{Key Categories}

% Category 1:

% Category 2:

\subsubsection{Implementation Example}

\subsubsection{Interaction with Other Functions}

%---------------------------------------

\subsection{Protect}

\subsubsection{Primary Purpose}

\subsubsection{Key Categories}

% Category 1:

% Category 2:

\subsubsection{Implementation Example}

\subsubsection{Interaction with Other Functions}

%---------------------------------------

\subsection{Detect}

\subsubsection{Primary Purpose}

\subsubsection{Key Categories}

% Category 1:

% Category 2:

\subsubsection{Implementation Example}

\subsubsection{Interaction with Other Functions}

%---------------------------------------

\subsection{Respond}

\subsubsection{Primary Purpose}

\subsubsection{Key Categories}

% Category 1:

% Category 2:

\subsubsection{Implementation Example}

\subsubsection{Interaction with Other Functions}

%---------------------------------------

\subsection{Recover}

\subsubsection{Primary Purpose}

\subsubsection{Key Categories}

% Category 1:

% Category 2:

\subsubsection{Implementation Example}

\subsubsection{Interaction with Other Functions}

%%%%%%%%%%%%%%%%%%%%%%%%%%%%%%%%%%%%%%%%%%%%%%%%%%%%%

\section{Implementation Tiers}

%====================================================

\subsection{Tier 1: Partial}

%---------------------------------------

\subsection{Tier 2: Risk Informed}

%---------------------------------------

\subsection{Tier 3: Repeatable}

%---------------------------------------

\subsection{Tier 4: Adaptive}

%---------------------------------------

\subsection{Startup Tier Progression Roadmap}

% Scenario: Small e-commerce startup progressing from Tier 1 to Tier 3

\subsubsection{Tier 1 to Tier 2 Transition}

% Steps:

% Resource requirements:

% Timeline:

\subsubsection{Tier 2 to Tier 3 Transition}

% Steps:

% Resource requirements:

% Timeline:

% Roadmap Table goes here (see templates/roadmap.tex)

%---------------------------------------

\subsection{Why Tier 4 Is Not Always the Right Target}

%%%%%%%%%%%%%%%%%%%%%%%%%%%%%%%%%%%%%%%%%%%%%%%%%%%%%

\section{Profiles and Gap Analysis}

%====================================================

% Scenario: Fictional university --- 10,000 students, 2,000 staff

\subsection{Current Profile}

% Identifies existing cybersecurity capabilities across at least 3 core functions

% Current Profile Table goes here

%---------------------------------------

\subsection{Target Profile}

% Desired cybersecurity posture across the same functions

% Target Profile Table goes here

%---------------------------------------

\subsection{Gap Analysis}

% Minimum 5 gaps identified between Current and Target Profiles

% Gap Analysis Table goes here (see templates/gap-analysis.tex)

%---------------------------------------

\subsection{Prioritised Remediation Plan}

% Timelines and responsible parties for each identified gap

%%%%%%%%%%%%%%%%%%%%%%%%%%%%%%%%%%%%%%%%%%%%%%%%%%%%%

\section{Incident Response: Hospital Ransomware Attack}

%====================================================

% Scenario: Ransomware attack affecting patient records across 3 hospitals

\subsection{Respond and Recover Functions Mapping}

\subsubsection{Respond}

% Map detailed incident response activities to the Respond function

\subsubsection{Recover}

% Map detailed recovery activities to the Recover function

%---------------------------------------

\subsection{Role of Protect and Detect in Prevention}

% Explain how stronger implementation of Protect and Detect could have
% prevented or minimised the impact of the ransomware attack

%---------------------------------------

\subsection{Post-Incident Improvement Plan}

% Using the Govern and Identify functions to prevent future occurrences

\subsubsection{Govern}

\subsubsection{Identify}

%---------------------------------------

\subsection{Communication Strategy}

\subsubsection{Patients}

\subsubsection{Regulators}

\subsubsection{The Public}

%%%%%%%%%%%%%%%%%%%%%%%%%%%%%%%%%%%%%%%%%%%%%%%%%%%%%

\section*{Chapter Summary}

Summarize the key findings from this chapter and briefly transition to the COBIT 2019 IT Governance Framework discussed in the next chapter.
